\begin{center}
    \Large{\textbf{Аннотация}}
\end{center}

\noindent
В работе исследуется задача порождающего моделирования условного распределения медицинских изображений: построение измеримого отображения~$T : \mathbb{R}^d \to \mathbb{R}^d$, переводящего распределение нативных компьютерных томограмм~$\pi_N$ в распределение их контрастных аналогов~$\pi_C$ (и обратно). Формально требуется решить транспортную задачу $T_{\#}\pi_0 = \pi_1$ в~$d=512^2$-мерном пространстве при наличии парных наблюдений $(\mathbf{x}_N,\mathbf{x}_C)\sim\gamma$, $\gamma\in\Pi(\pi_N,\pi_C)$.

Предлагается метод \textbf{Medical Flow Matching~(MFM)}~--- условный флоу матчинг с линейным интерполянтом $\mathbf{x}_t=(1-t)\mathbf{x}_0+t\mathbf{x}_1$. Минимизируется симуляционно-свободный функционал
\[
    \mathcal{L}_{\mathrm{CFM}}(\theta) = \mathbb{E}_{t\sim\mathcal{U}[0,1],\,(\mathbf{x}_0,\mathbf{x}_1)\sim\gamma}\bigl[\lVert f_\theta(\mathbf{x}_t,t)-(\mathbf{x}_1-\mathbf{x}_0)\rVert_2^2\bigr],
\]
аппроксимирующий маргинальное поле скоростей~$u_t$, порождающее перенос. Реализацией~$f_\theta$ служит архитектура \textbf{TimeResNet}~--- U-Net с поблочной аффинной модуляцией синусоидального временного эмбеддинга и свёрточным модулем самовнимания на боттлнеке. В отличие от диффузионных моделей, требующих $T\!\sim\!10^2$--$10^3$ шагов денойзинга, при линейном интерполянте кинетическая энергия потока минимальна, кривизна траектории нулевая, и интегрирование ОДУ методом Эйлера за один шаг точно с точностью до аппроксимационной ошибки сети.

\smallskip
\noindent\textit{Численные результаты на отложенной выборке из 20 пациентов (9\,372 среза):}

\smallskip
\noindent
\begin{tabular}{@{}lccc@{}}
\toprule
\textbf{Метрика} & \textbf{MFM} & \textbf{SMILE}~\cite{liu2025see} & \textbf{DDPM} \\
\midrule
MAE, HU $\downarrow$              & $\mathbf{4.25\pm1.09}$  & $30.22\pm4.95$ & $18.20\pm1.39$ \\
SSIM $\uparrow$                    & $\mathbf{0.994\pm0.007}$ & $0.887\pm0.026$ & $0.934\pm0.008$ \\
PSNR, дБ $\uparrow$               & $\mathbf{42.19\pm1.67}$ & $26.70\pm1.33$ & $30.05\pm1.48$ \\
Время на срез, с $\downarrow$     & $\mathbf{0.125}$        & $9.03$         & $44.20$ \\
\bottomrule
\end{tabular}

\smallskip
\noindent
Таким образом, MFM сокращает время инференса в \textbf{354~раза} по сравнению с DDPM и в \textbf{7.1~раз} снижает MAE по сравнению со SMILE. Из симметрии линейного интерполянта следует двунаправленность: единственная сеть~$f_\theta$ реализует трансляцию в обе стороны без увеличения числа параметров. Прикладная значимость подтверждается тем, что nnUNet\,v2, обученный на синтезированных MFM нативных срезах, достигает \textbf{Dice~=~0.926} (95.9\% от уровня 0.966 на реальных данных) в задаче сегментации аорты.

\medskip
\noindent\textbf{Ключевые слова:} условный флоу матчинг, оптимальный транспорт, обыкновенные дифференциальные уравнения, глубокое обучение, U-Net, трансляция КТ-изображений, синтез медицинских изображений.
